\documentclass{situationist-preprint}

\title{Theism and Additional Meaning}
\msid{EXI-11-b}
\msfamily{EXI-11}
\msdiscipline{Meaning, Death and the Absurd}
\mstype{Companion manuscript}
\mskeywords{meaning in life; theism; anti-theism; degrees of meaning; Kahane}

\begin{document}
\maketitle

\begin{abstract}
EXI-11-a argues that meaning does not require a cosmic purpose, and that a theistic framework, if true, would nevertheless make more meaning available (Proposition 4.1). The second claim is challenged from two sides. The supernaturalist says that what God adds is not more of the same but meaning as such; the anti-theist, following Kahane, says that God's existence would take meaning away, by removing the privacy, independence and self-direction on which some meaningful lives depend. We argue that the additional meaning is a difference of degree, and that the anti-theist's losses concern particular life-plans, not the objects of worth available to engagement, so that Proposition 4.1 survives both challenges.
\end{abstract}

\section{What theism adds}

EXI-11-a defines meaning as active engagement with what is worthwhile (Definition~2.1) and holds that meaning grows with the worth of what one engages and with the extent to which engagement relates one to more of what is worthwhile (Principle~P2; on degrees of meaning, see \citealp{Metz2013}). It identifies three additions a theistic framework would make: a being of maximal worth with whom one could stand in relation, a purpose unifying one's projects, and the preservation of what is good. Since the first and third are objects and relations of engagement of great worth that a naturalistic world lacks, Proposition~4.1 of EXI-11-a follows from Principle~P2.

Two points about the claim should be explicit. It is conditional: it concerns what a theistic framework would add if true, and so it holds whatever the answer to REL-01. And it concerns \emph{availability}: it says what a life could engage with, not that every life in a theistic world is more meaningful than every life in a naturalistic one.

\begin{proposition}\label{p:degree}
The meaning that a theistic framework would add is of the same kind as the meaning available without it, and differs in degree.
\end{proposition}

\begin{proof}
By Definition~2.1 of EXI-11-a, any meaning is engagement with what is worthwhile. Engagement with God, or with a purpose that unifies one's projects, is engagement with what is worthwhile; so the meaning it yields falls under the same definition. By Theorem~3.3 of EXI-11-a, lives without it can be meaningful, so it is not meaning as such. By Principle~P2 of EXI-11-a, it can increase the degree of meaning.
\end{proof}

\section{Two challenges}

\begin{objection}[Meaning as such]
\citet{Cottingham2003} holds that without God our projects lack an ultimate point; what theism adds is not more meaning but the only meaning that would not be undercut.
\end{objection}

\begin{reply}
The claim that naturalistic meaning is undercut is the claim that it depends on a cosmic purpose, which Theorem~3.3 of EXI-11-a rejects. The worth of knowledge, of persons and of beauty does not depend on a role assigned by a designer (Lemma~3.2 of EXI-11-a). What remains true in the objection is that a theistic framework would offer a single good towards which all one's projects could be ordered; Proposition~\ref{p:degree} counts that as an addition of degree.
\end{reply}

\begin{objection}[Anti-theism]
\citet{Kahane2011} argues that a person may reasonably prefer that God not exist, since God's existence would deprive her of privacy, of independence and of the possibility of understanding the world through her own inquiry. For such a person a theistic world would contain less meaning.
\end{objection}

\begin{reply}
We grant that the goods Kahane names are real and that some life-plans depend on them. But they bear on Proposition~4.1 of EXI-11-a in a different way from the goods theism adds. The worth of knowledge, persons and beauty is not diminished by God's existence, so the objects of engagement available without God remain available with God. What would be lost is a particular manner of engaging, unobserved and unaided, which is a component of some lives. The proposition concerns the meaning made available by a framework, and the framework adds objects of greater worth (Principle~P2 of EXI-11-a) while leaving the others in place. That a given person might reasonably prefer the world without them is compatible with this, as a person may reasonably prefer a smaller city whose ways she knows to a larger one with more to offer. Kahane's argument shows that the addition has costs for some lives; it does not show that the framework makes less meaning available.
\end{reply}

\section{Conclusion}

Theism, if true, would add to the meaning available in a life by adding objects and relations of engagement of great worth. The addition is one of degree, which is why its absence leaves meaning in place, and it is compatible with the costs to some lives that the anti-theist identifies.

\begin{thebibliography}{9}\small
\bibitem[Cottingham(2003)]{Cottingham2003} Cottingham, J. (2003). \emph{On the Meaning of Life}. Routledge.
\bibitem[Kahane(2011)]{Kahane2011} Kahane, G. (2011). Should we want God to exist? \emph{Philosophy and Phenomenological Research}, 82(3), 674--696.
\bibitem[Metz(2013)]{Metz2013} Metz, T. (2013). \emph{Meaning in Life: An Analytic Study}. Oxford University Press.
\end{thebibliography}

\end{document}
